\fancyhead[L]{\small NOUETS BEOBACHTUNGEN IN ÄGYPTEN.}
\hypertarget{native-539Nouet}{}\label{native-539Nouet}
\begin{center}\small Göttingische gelehrte Anzeigen. Stück 79. S. 780..781. 18. Mai 1811.\\[3pt]\rule{.72\linewidth}{.3pt}\end{center}

\emph{Beschreibung Ägyptens. Erster Band, gedruckt in Paris bei der Kaiserlichen Druckerei} 1809. \raisebox{-0.3ex}{,}\emph{Erste Abhandlung. Astronomische Beobachtungen in Ägypten in den Jahren} VI. VII \emph{und} VIII (1798. 1799. 1800) \emph{von} M. \textsc{Nouet} \emph{Astronom der Kommission für Wissenschaften und Künste Ägyptens.}

\textsc{Nouet} verwendete einen Multiplikationskreis von 25 Zentimetern Durchmesser mit Dezimaleinteilung (weiter unten wird auch ein Quadrant mit 35 Zentimetern Radius erwähnt), eine Seeuhr von \textsc{Louis Berthoud} und ein achromatisches Fernrohr von \textsc{Dollond} mit 63 Millimetern Öffnung. In Kairo erhielt er von \textsc{Beauchamp}, den er dort traf, noch einen zweiten Chronometer. Die Beobachtungen ergaben die Breiten und Längen von 36 verschiedenen Punkten. Der Platz und der Zweck unserer Zeitschrift erlauben es nicht, näher auf die Einzelheiten einzugehen. Wir begnügen uns damit, nur einiges daraus hervorzuheben. Die südlichsten Punkte sind die Insel \emph{Philä} oberhalb der Katarakte des Nils auf $24^\circ\,1'\,34''$ Breite und \emph{Syene} auf $24^\circ\,5'\,23''$ Breite. Von diesem letzten Ort kennt man die Sage, dass zur Zeit der Sommersonnenwende ein Brunnen bis auf den Grund erleuchtet worden sei. \textsc{Nouet} berechnet nach \textsc{Laplace}s Formel, wann die Schiefe der Ekliptik jener Breite gleich war, und setzt diesen Zeitpunkt 3430 Jahre vor unserer Zeitrechnung an. Deshalb glaubt er, die Sage so weit zurückdatieren zu müssen. Dieser Schluss scheint jedoch nicht hinreichend begründet, da ein solches Phänomen auch dann noch eintreten konnte, wenn der \emph{Mittelpunkt} der Sonne in geringem Abstand südlich am Zenit vorbeiging. Die Längen wurden größtenteils chronometrisch bestimmt. \textsc{Nouet} hatte aber auch Gelegenheit, vier Bedeckungen zu beobachten. Er beobachtete die Bedeckung der Venus am 13. Dezember 1798 in Ssalehhiyeh, die Bedeckung von $\delta$ im Skorpion am 21. April 1799 in Kairo, die Bedeckung der Venus am 23. November 1799 am selben Ort und die Bedeckung von $\alpha$ im Skorpion am 28. Juli 1800 in Alexandria. Außerdem beobachtete er in Kairo eine beträchtliche Anzahl von Verfinsterungen der Jupitermonde. Die Abweichung der Magnetnadel in Alexandria im Mittel aus 26 Beobachtungen, deren Datum jedoch nicht angegeben ist. ergab $13^\circ\,6'$ N.\,W., und die Neigung im Mittel aus 12 Beobachtungen $47^\circ\,30'$ nördlich.
